\section{Conclusion and Future Work}\label{sec:conclusions-and-future-work}

This work delivers a pipeline that combines authorisation data from different sources and a BCA ranker for crop
and pathogen queries. The system combines support counts with context similarity, keeps products as evidence, and
reports known and inferred candidates separately. The evaluation shows retrieval of observed authorisation
associations, not biological efficacy. The research question is therefore answered only for warm-start retrieval of
authorisation-derived associations: the system can recover held-out associations in the current Portuguese catalogue,
but it cannot establish efficacy, approval, or suitability for new uses.

Future work should first add structured sources with verified crop--pathogen pairs and evaluate the ranker with
independent efficacy or field data. It should then use a validation split to tune weights and test cold-start and
temporal settings. A bounded large language model (LLM) stage may assist irregular source mapping, but deterministic
validation and human review must check every suggestion before it enters the modelling table. Inferred mappings
cannot count as verified authorisation evidence.